\chapter{Formulación matemática del problema}
\label{ch:formulacion}

Este capítulo es el núcleo matemático del trabajo. Se formula
el modelo de selección de carteras de Markowitz como un problema de
optimización cuadrática convexa (QP), se analizan sus condiciones de
optimalidad de Karush-Kuhn-Tucker (KKT), se estudia la dualidad y se
presenta la solución analítica de la frontera eficiente para el caso sin
restricciones de desigualdad (cuya derivación completa se recoge en el
Apéndice~\ref{ap:frontera-analitica}). La notación y los resultados aquí
desarrollados son la base sobre la que se construye la implementación
computacional del Capítulo~\ref{ch:metodologia}.

%--------------------------------------------------------------------
\section{El modelo media-varianza como problema de optimización}
%--------------------------------------------------------------------

\subsection{Notación}

Sea un universo de $n$ activos financieros. Se denota:

\begin{itemize}
    \item $\mathbf{w} = (w_1, \dots, w_n)^T \in \mathbb{R}^n$: vector
          de pesos de la cartera.
    \item $\boldsymbol{\mu} = (\mu_1, \dots, \mu_n)^T \in \mathbb{R}^n$:
          vector de rentabilidades esperadas, donde $\mu_i = \E[r_i]$.
    \item $\boldsymbol{\Sigma} = [\sigma_{ij}] \in \mathbb{R}^{n \times n}$:
          matriz de covarianza, simétrica y semidefinida positiva:
          $\boldsymbol{\Sigma} \succeq 0$.
    \item $\mathbf{1} = (1, \dots, 1)^T \in \mathbb{R}^n$.
\end{itemize}

\subsection{Formulación primal}

El problema de Markowitz en su forma estándar (minimizar riesgo sujeto a
una rentabilidad objetivo) se escribe:

\begin{equation}
    \boxed{
    \begin{aligned}
        \min_{\mathbf{w} \in \mathbb{R}^n} \quad
        & f(\mathbf{w}) = \frac{1}{2} \mathbf{w}^T \boldsymbol{\Sigma}
          \mathbf{w} \\
        \text{sujeto a} \quad
        & \mathbf{w}^T \boldsymbol{\mu} \ge \mu^*, \\
        & \mathbf{1}^T \mathbf{w} = 1, \\
        & w_i \ge 0, \quad i = 1, \dots, n, \\
        & w_i \le w_{\max}, \quad i = 1, \dots, n.
    \end{aligned}}
    \label{eq:qp-completo}
\end{equation}

El factor $\frac{1}{2}$ en la función objetivo es puramente algebraico:
simplifica las expresiones de las derivadas sin alterar el óptimo.
La primera restricción exige que la cartera alcance al menos una
rentabilidad esperada $\mu^*$. La segunda es la restricción de
presupuesto. Las dos últimas son restricciones de desigualdad que
acotan inferior y superiormente cada peso.

El caso particular en que se omite la restricción de rentabilidad objetivo
($\mathbf{w}^T \boldsymbol{\mu} \ge \mu^*$) se conoce como \textbf{cartera
de mínima varianza global} (GMV):
\begin{equation}
    \min_{\mathbf{w}} \;
    \frac{1}{2} \mathbf{w}^T \boldsymbol{\Sigma} \mathbf{w}
    \quad \text{s.a.} \quad
    \mathbf{1}^T \mathbf{w} = 1, \;
    \mathbf{w} \ge \mathbf{0}, \;
    \mathbf{w} \le w_{\max} \mathbf{1}.
    \label{eq:gmv-formal}
\end{equation}

%--------------------------------------------------------------------
\section{Convexidad del problema cuadrático}
%--------------------------------------------------------------------

Un problema de optimización es \textbf{convexo} si la función objetivo es
convexa y la región factible es un conjunto convexo. Esta propiedad es crucial:
garantiza que todo mínimo local es también global y que las condiciones KKT son
necesarias y suficientes para la optimalidad (junto con la condición de Slater,
Sección~\ref{sec:dualidad}) \cite{Boyd2004}. A continuación se
comprueba que el problema de Markowitz cumple ambas condiciones.

\subsection{Convexidad de la función objetivo}

La función objetivo es la forma cuadrática:
\begin{equation}
    f(\mathbf{w}) = \frac{1}{2} \mathbf{w}^T \boldsymbol{\Sigma} \mathbf{w}.
    \label{eq:funcion-objetivo}
\end{equation}

Su matriz hessiana es $\nabla^2 f(\mathbf{w}) = \boldsymbol{\Sigma}$. Dado
que $\boldsymbol{\Sigma}$ es una matriz de covarianza, es semidefinida
positiva por construcción:
\begin{equation}
    \mathbf{x}^T \boldsymbol{\Sigma} \mathbf{x}
    = \Var[\mathbf{x}^T \mathbf{r}] \ge 0,
    \quad \forall \mathbf{x} \in \mathbb{R}^n.
    \label{eq:psd}
\end{equation}

Por tanto, $f$ es una función convexa ($\boldsymbol{\Sigma} \succeq 0$
implica convexidad; si $\boldsymbol{\Sigma} \succ 0$, la función es
estrictamente convexa). En la práctica, con $T > n$ observaciones y
activos no redundantes, $\boldsymbol{\Sigma}$ es definida positiva y la
solución es única.

\subsection{Convexidad de la región factible}

La región factible $\mathcal{F}$ está definida por:
\begin{equation}
    \mathcal{F} = \left\{
        \mathbf{w} \in \mathbb{R}^n \;\middle|\;
        \begin{aligned}
            &\mathbf{1}^T \mathbf{w} = 1, \\
            &\mathbf{w}^T \boldsymbol{\mu} \ge \mu^*, \\
            &0 \le w_i \le w_{\max}, \; i = 1, \dots, n
        \end{aligned}
    \right\}.
    \label{eq:region-factible}
\end{equation}

$\mathcal{F}$ es intersección de conjuntos convexos ---el hiperplano
$\mathbf{1}^T \mathbf{w} = 1$, el semiespacio $\mathbf{w}^T \boldsymbol{\mu} \ge
\mu^*$ y el hiperrectángulo $0 \le w_i \le w_{\max}$---, luego es convexa; y es
compacta (cerrada y acotada), de modo que el teorema de Weierstrass garantiza la
existencia de un mínimo global. Al ser convexos tanto el objetivo como la
región, el problema~(\ref{eq:qp-completo}) es un \textbf{programa cuadrático
convexo} (QP), y el óptimo que halla un solver estándar (\texttt{cvxpy}) es el
óptimo global.

Geométricamente, $\mathcal{F}$ es un politopo convexo: la intersección del
$(n-1)$-simplex $\{\mathbf{w} \ge \mathbf{0} \mid \mathbf{1}^T \mathbf{w} = 1\}$
con el hipercubo $[0, w_{\max}]^n$ y el semiespacio $\mathbf{w}^T
\boldsymbol{\mu} \ge \mu^*$. Los topes $w_i \le w_{\max}$ truncan los vértices
del simplex ---las carteras concentradas en un solo activo---, lo que hace la
región más realista pero impide una solución analítica cerrada; y si $\mu^*$ es
demasiado grande, $\mathcal{F}$ puede quedar vacía y el problema resulta
infactible.

%--------------------------------------------------------------------
\section{Condiciones de Karush-Kuhn-Tucker (KKT)}
%--------------------------------------------------------------------

Las condiciones KKT caracterizan el óptimo de un problema con restricciones de
igualdad y desigualdad a través de su Lagrangiano; su forma general se recuerda
en el Apéndice~\ref{ap:kkt-general}. A continuación se particularizan al QP de
Markowitz.

Escribiendo el QP~(\ref{eq:qp-completo}) en la forma estándar general ---con
función objetivo $f(\mathbf{w}) = \frac{1}{2}\mathbf{w}^T \boldsymbol{\Sigma}
\mathbf{w}$, las desigualdades $g_i \le 0$ (long-only, tope por activo y
rentabilidad objetivo) y la igualdad $\mathbf{1}^T\mathbf{w} = 1$--- e
introduciendo los multiplicadores $\alpha_i, \beta_i \ge 0$ (cotas por activo),
$\gamma \ge 0$ (rentabilidad) y $\lambda \in \mathbb{R}$ (presupuesto), el
Lagrangiano resultante (detallado en el Apéndice~\ref{ap:kkt-markowitz}) da lugar
a las siguientes \textbf{condiciones KKT} para el óptimo $\mathbf{w}^*$:

\subsubsection*{1. Estacionariedad}

El gradiente del Lagrangiano respecto a $\mathbf{w}$ se anula en el óptimo:
\begin{equation}
    \nabla_{\mathbf{w}} \mathcal{L}
    = \boldsymbol{\Sigma} \mathbf{w}^*
    - \boldsymbol{\alpha} + \boldsymbol{\beta}
    - \gamma \boldsymbol{\mu}
    - \lambda \mathbf{1}
    = \mathbf{0}.
    \label{eq:kkt-estacionariedad}
\end{equation}

Esta ecuación tiene una lectura financiera directa: en el óptimo, el riesgo
adicional que aporta cada activo ($[\boldsymbol{\Sigma} \mathbf{w}^*]_i$) se
equilibra con su beneficio en rentabilidad ($\gamma \mu_i$), ajustado por los
multiplicadores de las restricciones activas y el precio del presupuesto
$\lambda$.

\subsubsection*{2. Factibilidad primal y dual}

El punto $\mathbf{w}^*$ satisface todas las restricciones y los multiplicadores
de las desigualdades son no negativos:
\begin{equation}
    \mathbf{1}^T \mathbf{w}^* = 1, \;\;
    0 \le w_i^* \le w_{\max}, \;\;
    \mathbf{w}^{*T} \boldsymbol{\mu} \ge \mu^*;
    \qquad
    \alpha_i, \beta_i, \gamma \ge 0.
    \label{eq:kkt-primal}
\end{equation}

\subsubsection*{3. Holgura complementaria}

Para cada restricción de desigualdad, o bien la restricción está activa
(se cumple con igualdad) o bien su multiplicador es cero:
\begin{align}
    \alpha_i \, w_i^* &= 0, \quad i = 1, \dots, n,
        \label{eq:kkt-holgura1} \\
    \beta_i \, (w_i^* - w_{\max}) &= 0, \quad i = 1, \dots, n,
        \label{eq:kkt-holgura2} \\
    \gamma \, (\mu^* - \mathbf{w}^{*T} \boldsymbol{\mu}) &= 0.
        \label{eq:kkt-holgura3}
\end{align}

La holgura complementaria admite una interpretación económica clara según el
peso del activo en el óptimo: si $0 < w_i^* < w_{\max}$, entonces
$\alpha_i = \beta_i = 0$ y la estacionariedad se reduce a
$[\boldsymbol{\Sigma}\mathbf{w}^*]_i = \gamma\mu_i + \lambda$ (el coste marginal
de riesgo iguala al beneficio marginal de rentabilidad); si $w_i^* = 0$, el
activo queda fuera porque su contribución al riesgo supera a la de rentabilidad;
y si $w_i^* = w_{\max}$, está en su tope porque su rentabilidad justifica asumir
más riesgo del que permiten las restricciones.

La implementación verifica explícitamente que la solución de \texttt{cvxpy}
satisface estas condiciones: el test \texttt{test\_condiciones\_kkt} recupera
los multiplicadores a partir de la solución numérica y comprueba las cuatro
condiciones con tolerancia $10^{-4}$, una validación independiente del solver.

%--------------------------------------------------------------------
\section{Dualidad e interpretación económica}
\label{sec:dualidad}
%--------------------------------------------------------------------

\subsection{El problema dual}

Asociado al problema primal (\ref{eq:qp-completo}) existe un
\textbf{problema dual} que proporciona una cota inferior del valor
óptimo primal. Si el problema convexo satisface la condición de Slater (existe un punto
estrictamente factible), la \textbf{dualidad fuerte} garantiza que los valores
óptimos primal y dual coinciden \cite{Boyd2004}.

En nuestro QP, Slater se cumple si $n \cdot w_{\max} > 1$ (los topes son
compatibles con el presupuesto) y $\mu^*$ es estrictamente menor que la máxima
rentabilidad esperada alcanzable respetando los topes; bajo esas condiciones
existe un punto estrictamente interior a la región factible.

\subsection{Interpretación de los multiplicadores}

Los multiplicadores de KKT se interpretan como \textbf{precios sombra}: el
precio sombra de una restricción indica cuánto bajaría la varianza mínima de la
cartera si esa restricción se relajara en una unidad. Dicho de otro modo, mide
lo que ``cuesta'' esa restricción en el óptimo. En nuestro problema:
\begin{itemize}
    \item $\lambda$ (presupuesto): cuánto cambiaría el riesgo si el capital a
          invertir fuera algo mayor o menor.
    \item $\gamma$ (rentabilidad objetivo): cuánto bajaría el riesgo si se
          exigiera un poco menos de rentabilidad.
    \item $\alpha_i$ (cota $w_i \ge 0$): si es positivo, permitir una pequeña
          posición corta en el activo $i$ reduciría el riesgo.
    \item $\beta_i$ (tope $w_i \le w_{\max}$): si es positivo, dejar invertir
          algo más del tope en el activo $i$ reduciría el riesgo.
\end{itemize}

%--------------------------------------------------------------------
\section{Solución analítica de la frontera eficiente}
%--------------------------------------------------------------------

Cuando se eliminan las restricciones de desigualdad ($w_i \ge 0$ y
$w_i \le w_{\max}$) y se supone $\boldsymbol{\Sigma} \succ 0$, el problema
admite una solución analítica cerrada, uno de los resultados clásicos de la
teoría de carteras \cite{Merton1972}. La frontera de mínima varianza es una
hipérbola en el plano $(\sigma_p, \mu_p)$,
\begin{equation*}
    \sigma_p^2 = \frac{A \mu_p^2 - 2B \mu_p + C}{\Delta},
    \qquad
    A = \mathbf{1}^T \boldsymbol{\Sigma}^{-1} \mathbf{1}, \;
    B = \mathbf{1}^T \boldsymbol{\Sigma}^{-1} \boldsymbol{\mu}, \;
    C = \boldsymbol{\mu}^T \boldsymbol{\Sigma}^{-1} \boldsymbol{\mu},
\end{equation*}
con $\Delta = AC - B^2 > 0$ (positivo por Cauchy-Schwarz siempre que
$\boldsymbol{\mu}$ y $\mathbf{1}$ sean linealmente independientes). Su vértice
es la cartera de mínima varianza global (GMV), con $\mu_{\text{GMV}} = B/A$, y
el punto de tangencia con una recta por el origen da la cartera de máximo ratio
de Sharpe:
\begin{equation*}
    \mathbf{w}_{\text{GMV}}
    = \frac{\boldsymbol{\Sigma}^{-1} \mathbf{1}}
           {\mathbf{1}^T \boldsymbol{\Sigma}^{-1} \mathbf{1}},
    \qquad
    \mathbf{w}_{T}
    = \frac{\boldsymbol{\Sigma}^{-1} \boldsymbol{\mu}}
           {\mathbf{1}^T \boldsymbol{\Sigma}^{-1} \boldsymbol{\mu}}.
\end{equation*}
La derivación completa por multiplicadores de Lagrange ---con la ecuación de la
frontera, las expresiones de la GMV y de la cartera tangente, y su validación
numérica frente a \texttt{cvxpy}--- se recoge en el
Apéndice~\ref{ap:frontera-analitica}.

%--------------------------------------------------------------------
\section{Del caso analítico al caso con restricciones}
%--------------------------------------------------------------------

Cuando se añaden las restricciones $w_i \ge 0$ y $w_i \le w_{\max}$,
la solución deja de admitir una expresión cerrada. La frontera eficiente
ya no es una hipérbola perfecta, sino una curva suave definida a trozos
por los puntos donde las restricciones se activan o desactivan, y la
resolución pasa a ser numérica.

Sin embargo, las propiedades cualitativas se conservan:
\begin{itemize}
    \item El problema sigue siendo convexo (la región factible es un
          subconjunto convexo del simplex).
    \item La solución es única si $\boldsymbol{\Sigma} \succ 0$.
    \item Las condiciones KKT siguen caracterizando completamente el
          óptimo.
    \item La frontera eficiente es creciente y cóncava en el espacio
          $(\sigma_p, \mu_p)$.
\end{itemize}

Esto justifica el uso de \texttt{cvxpy} como solver exacto: la convexidad
garantiza optimalidad global, y las condiciones KKT proporcionan una
herramienta de verificación independiente del solver.
